% Part: proof-theory
% Chapter: natural-deduction
% Section: rules-proofs

\documentclass[../../../include/open-logic-section]{subfiles}

\begin{document}

\olfileid[id]{pt}{nat}{rp}

\olsection{Aturan dan \usetoken{P}{proof}}

Kita mulai dengan memberikan beberapa definisi dan menetapkan sejumlah
terminologi.

Tinjau dua aturan dalam \Log{N1c} dan~\Log{N1i} yang menyebut
$!A \lif !B$.
\[
\bottomAlignProof
\AxiomC{$\Discharge{!A}{x}$}
\shortDeduce
\DeduceC{$!B$}
\DischargeRule{\Intro{\lif}}{x}
\UnaryInfC{$!A \lif !B$}
\DisplayProof
\qquad
\bottomAlignProof
\AxiomC{$!A \lif !B$}
\AxiomC{$!A$}
\RightLabel{\Elim{\lif}}
\BinaryInfC{$!B$}
\DisplayProof
\]
Aturan \Elim{\lif} cukup sederhana. !!{formula} di atas garis disebut
\emph{premis}. !!{formula} di sebelah kiri adalah !!{formula}
\emph{utama}~$!A \lif !B$. Semua aturan !!{elimination} memiliki satu
premis semacam itu, selalu di posisi paling kiri. Premis itu disebut
\emph{premis mayor} inferensi. Jika ada premis lain, seperti dalam kasus
ini, premis-premis itu disebut \emph{premis minor}. Premis bagi aturan
!!{introduction} juga disebut premis mayor. !!{formula} di bawah garis
disebut \emph{kesimpulan}.

Aturan \Intro{\lif} hanya memiliki satu premis, dan di sini kesimpulannya
adalah formula utama~$!A \lif !B$. Semua aturan !!{introduction} memiliki
bentuk seperti ini: kesimpulannya adalah formula utama, dan operator
utamanya merupakan operator yang ``diintroduksi.''

Notasi $\Discharge{!A}{x}$ bermakna sebagai berikut. Dalam !!a{proof},
inferensi \Intro{\lif} diterapkan pada !!a{formula}~$!B$ sebagai premis.
Sub-!!{proof} yang berakhir pada~$!B$ mempunyai sejumlah !!{formula}
awal. Formula-formula itu disebut \emph{asumsi}. !!^a{proof} atas~$!B$
dari satu atau lebih asumsi berbentuk~$!A$ adalah !!a{proof} bahwa $!B$
mengikuti dari~$!A$. Ketika menerapkan aturan \Intro{\lif}, kita
menyimpulkan~$!A \lif !B$, dan kesimpulan itu tidak lagi bergantung pada
asumsi~$!A$. Untuk menunjukkan hal ini, asumsi-asumsi berbentuk~$!A$
dalam !!{proof} yang mencakup inferensi~\Intro{\lif} tersebut
\emph{dilepaskan} atau \emph{ditutup}. Aturannya menentukan asumsi mana
yang dapat dilepaskan. Dalam kasus \Intro{\lif} dengan kesimpulan
$!A \lif !B$, asumsi itu berbentuk~$!A$, yakni sama dengan anteseden
kondisional. Kadang-kadang penting untuk menunjukkan asumsi mana yang
dilepaskan pada aturan mana, dan hal ini dilakukan dengan label
pelepasan~$x$. Yang penting, aturan yang melepaskan asumsi tidak
\emph{harus} melepaskan \emph{semua} asumsi dengan bentuk yang
disyaratkan, bahkan tidak harus melepaskan asumsi apa pun. Misalnya,
berikut adalah konstruksi yang benar, meskipun \Intro{\lif} pertama tidak
melepaskan asumsi apa pun:
\[
\AxiomC{$\Discharge{!A}{y}$}
\DischargeRule{\Intro{\lif}}{x}
\UnaryInfC{$!B \lif !A$}
\DischargeRule{\Intro{\lif}}{y}
\UnaryInfC{$!A \lif (!B \lif !A)$}
\DisplayProof
\]
Ketika suatu aturan tidak melepaskan asumsi apa pun, kita cukup
menghilangkan label pelepasan (seperti yang dilakukan di sini).

Dalam !!{proof}
\[
\AxiomC{$\Discharge{!A}{x}$}
\AxiomC{$\Discharge{!A}{y}$}
\RightLabel{\Intro{\land}}
\BinaryInfC{$!A \land !A$}
\RightLabel{\Elim{\land}}
\UnaryInfC{$!A$}
\DischargeRule{\Intro{\lif}}{x}
\UnaryInfC{$!A \lif !A$}
\DischargeRule{\Intro{\lif}}{y}
\UnaryInfC{$!A \lif (!A \lif !A)$}
\DisplayProof
\]
inferensi \Intro{\lif} pertama melepaskan asumsi kiri $!A^x$, sedangkan
inferensi kedua melepaskan asumsi kanan~$!A^y$. Dengan kata lain, semua
asumsi berlabel~$x$ dilepaskan oleh \Intro{\lif} berlabel~$x$, dan semua
asumsi berlabel~$y$ dilepaskan oleh inferensi berlabel~$y$. Agar hal ini
berfungsi, semua asumsi dengan label yang sama juga harus berupa
!!{formula} yang sama.

Aturan kuantor sedikit lebih rumit. Misalnya, aturan untuk $\lexists$
dalam \Log{N1i} adalah:
\[
\bottomAlignProof
\AxiomC{$!A(t)$}
\RightLabel{\Intro{\lexists}}
\UnaryInfC{$\lexists[x][!A(x)]$}
\DisplayProof
\qquad
\bottomAlignProof
\AxiomC{$\lexists[x][!A(x)]$}
\AxiomC{$\Discharge{!A(c)}{x}$}
\shortDeduce
\DeduceC{$!B$}
\DischargeRule{\Elim{\lexists}}{x}
\BinaryInfC{$!B$} \DisplayProof
\bottomAlignProof
\]
Dalam aturan \Intro{\lexists}, premisnya ialah~$!A(t)$, dengan $t$
sebagai suku tertutup, yakni suku tanpa variabel. Inferensi semacam itu
benar---yakni cocok dengan aturan skematis \Intro{\lexists}---jika
terdapat suatu !!{formula}~$!A$, suatu variabel~$x$, dan suatu suku
tertutup~$t$ sedemikian sehingga kesimpulannya
$\lexists[x][!A]$ dan premisnya~$\Subst{!A}{t}{x}$.
Aturan \Elim{\lexists} memungkinkan kita melepaskan asumsi berbentuk
$!A(c)$; yakni, bagi setiap inferensi terdapat suatu !!{constant}~$c$,
dan asumsi berbentuk $\Subst{!A}{c}{x}$ dapat dilepaskan. Semua asumsi
yang dilepaskan oleh satu inferensi harus menggunakan~$c$ yang sama.
Konstanta ini disebut \emph{variabel eigen} dari inferensi
\Elim{\lexists}. Inferensi tersebut hanya benar jika $c$ tidak muncul
dalam asumsi terbuka selain $!A(c)$ yang dilepaskan, tidak muncul
dalam~$!B$, dan tidak muncul dalam~$!A$. Ketentuan ini disebut
\emph{syarat variabel eigen}.

!!^{proof} dalam \Log{N1c} dan~\Log{N1i} adalah pohon berhingga formula
yang dibentuk secara induktif dari aksioma dengan menggunakan
aturan inferensi. Setiap daun merupakan asumsi, yang mungkin diberi label
pelepasan, dan setiap !!{formula} lainnya mengikuti dari
semua !!{formula} tepat di atasnya melalui salah satu aturan
inferensi sistem. Jika suatu kemunculan !!{formula} sebagai asumsi dalam
pohon di atas suatu inferensi belum dilepaskan oleh inferensi di bawahnya,
kemunculan itu disebut \emph{terbuka} (pada inferensi tersebut).

\begin{defn}\ollabel{defn:proof}
\emph{!!^{proof}} dalam \Log{N1c} dan~\Log{N1i} adalah pohon berhingga
!!{formula} yang didefinisikan secara induktif sebagai berikut:
\begin{enumerate}
  \item\ollabel{defn:prf:assumption}
  Setiap !!{formula} yang diberi label pelepasan seperti $x$ atau $y$
  merupakan !!a{proof} dengan sendirinya. Dalam !!a{proof} yang hanya
  terdiri atas satu !!{formula}, !!{formula} tersebut dihitung sebagai
  asumsi terbuka.
  \item\ollabel{defn:prf:one-prem}
  Jika $R$ adalah aturan dengan satu, dua, atau tiga premis $!A_1$,
  $!A_2$, $!A_3$, kesimpulan~$!A$, dan $\delta_1$, $\delta_2$,
  $\delta_3$ masing-masing merupakan !!{proof} yang berakhir pada
  $!A_1$, $!A_2$, $!A_3$, maka
  \[
  \AxiomC{}
  \RightLabel{$\delta_1$}
  \DeduceC{$!A_1$}
  \RightLabel{$R$}
  \UnaryInfC{$!A$}
  \DisplayProof
\qquad
  \AxiomC{}
  \RightLabel{$\delta_1$}
  \DeduceC{$!A_1$}
  \AxiomC{}
  \RightLabel{$\delta_2$}
  \DeduceC{$!A_2$}
  \RightLabel{$R$}
  \BinaryInfC{$!A$}
  \DisplayProof
\qquad
  \AxiomC{}
  \RightLabel{$\delta_1$}
  \DeduceC{$!A_1$}
  \AxiomC{}
  \RightLabel{$\delta_2$}
  \DeduceC{$!A_2$}
  \AxiomC{}
  \RightLabel{$\delta_3$}
  \DeduceC{$!A_3$}
  \RightLabel{$R$}
  \TrinaryInfC{$!A$}
  \DisplayProof
  \]
  masing-masing merupakan !!a{proof}, asalkan label asumsi dalam
  !!{proof} tersebut kompatibel (tidak ada dua !!{formula}
  berbeda yang memiliki label pelepasan yang sama) dan
  inferensi-inferensinya merupakan penerapan~$R$ yang benar.
  \item Kemunculan !!{formula} paling atas dalam !!{proof} baru dihitung
  sebagai \emph{asumsi terbuka} dari !!{proof} tersebut jika kemunculan
  itu merupakan asumsi terbuka dari $\delta_1$, $\delta_2$,
  atau~$\delta_3$, kecuali kemunculan yang dilepaskan oleh aturan.
  Jika aturan diberi label~$x$ (atau $x\,y$), maka untuk
  \Intro{\lif} dan~\Intro{\lnot}, kemunculan yang dilepaskan ialah
  semua asumsi terbuka dalam~$\delta_1$ berlabel~$x$; untuk
  \Elim{\lexists}, semua asumsi terbuka berlabel~$x$
  dalam~$\delta_2$; dan untuk \Elim{\lor}, semua asumsi terbuka
  berlabel~$x$ dalam~$\delta_2$ serta semua asumsi terbuka berlabel~$y$
  dalam~$\delta_3$.
\end{enumerate}
\end{defn}

!!{proof} yang digunakan dalam konstruksi induktif suatu !!a{proof}
disebut \emph{sub-!!{proof}}-nya. Sub-!!{proof} yang berakhir pada premis
suatu inferensi disebut \emph{sub-!!{proof} langsung} dari inferensi
tersebut. Aturan \Log{N1c} dan \Log{N1i} diberikan dalam
\olref{tab:N1}.

\subfile{rules-N1}

\begin{defn}
\emph{!!{height}}~$\pheight{\delta}$ dari !!a{proof}~$\delta$
didefinisikan secara induktif sebagai berikut.
\begin{enumerate}
  \item Jika $\delta$ terdiri atas satu asumsi,
  $\pheight{\delta}=0$.
  \item Jika $\delta$ berakhir pada inferensi dengan satu premis dan
  sub-!!{proof} langsung~$\delta_1$, maka
  $\pheight{\delta}=\pheight{\delta_1}+1$.
  \item Jika $\delta$ berakhir pada inferensi dengan dua premis dan
  sub-!!{proof} langsung~$\delta_1$ dan~$\delta_2$, maka
  $\pheight{\delta}=
  \max(\pheight{\delta_1},\pheight{\delta_2})+1$.
  \item Jika $\delta$ berakhir pada \Elim{\lor} dengan
  sub-!!{proof} langsung~$\delta_1$, $\delta_2$, dan~$\delta_3$, maka
  $\pheight{\delta}=
  \max(\pheight{\delta_1},\pheight{\delta_2},\pheight{\delta_3})+1$.
\end{enumerate}
Jika suatu sekuen~$S$ memiliki !!a{proof} dengan tinggi $\le n$, kita
menulis $\Proves[n] S$.
\end{defn}

\begin{ex}
Dalam \Log{N1i}, setiap !!{formula} dihitung sebagai !!a{proof} dari
dirinya sendiri sebagai asumsi terbuka. Jadi,
\[
!A \land !B^x
\]
adalah !!a{proof}~$\delta_1$. Tingginya~$0$.
Dari $\delta_1$ kita dapat memperoleh dua !!{proof}~$\delta_2$
dan~$\delta_3$ dengan menerapkan salah satu dari dua versi
\Elim{\land}:
\[
\AxiomC{$!A \land !B^x$}
\RightLabel{\Elim{\land}[2]}
\UnaryInfC{$!B$}
\DisplayProof
\qquad
\AxiomC{$!A \land !B^x$}
\RightLabel{\Elim{\land}[1]}
\UnaryInfC{$!A$}
\DisplayProof
\]
Sub-!!{proof} langsung dari inferensi terakhir dalam~$\delta_2$ dan
dalam~$\delta_3$ sama, yaitu $!A \land !B$. Kedua !!{proof} tersebut
memiliki kemunculan $!A\land !B$ sebagai asumsi terbuka. Kita memperoleh
$\pheight{\delta_2}=\pheight{\delta_3}=1$.

Aturan \Intro{\land} memerlukan dua premis berbentuk $!B$ dan $!A$
untuk membenarkan kesimpulan berbentuk~$!B \land !A$. Karena itu,
berikut adalah !!a{proof}~$\delta_4$:
\[
\AxiomC{$!A \land !B^x$}
\RightLabel{\Elim{\land}[2]}
\UnaryInfC{$!B$}
\LeftSubproofLabel{$\delta_2$}
\AxiomC{$!A \land !B^x$}
\RightLabel{\Elim{\land}[1]}
\UnaryInfC{$!A$}
\RightSubproofLabel{$\delta_3$}
\RightLabel{\Intro{\land}}
\BinaryInfC{$!B \land !A$}
\DisplayProof
\]
Tingginya
$\pheight{\delta_4}=\max(\pheight{\delta_2},
\pheight{\delta_3})+1=\max(1,1)+1=2$. Bukti ini memiliki dua
kemunculan $!A \land !B$ sebagai asumsi, dan keduanya terbuka.

Akhirnya, kita dapat menerapkan \Intro{\lif} untuk memperoleh
$\delta_5$:
\[
\AxiomC{$\Discharge{!A \land !B}{x}$}
\RightLabel{\Elim{\land}[2]}
\UnaryInfC{$!B$}
\LeftSubproofLabel{$\delta_2$}
\AxiomC{$\Discharge{!A \land !B}{x}$}
\RightLabel{\Elim{\land}[1]}
\UnaryInfC{$!A$}
\RightSubproofLabel{$\delta_3$}
\RightLabel{\Intro{\land}}
\BinaryInfC{$!B \land !A$}
\DischargeRule{\Intro{\lif}}{x}
\UnaryInfC{$(!A \land !B) \lif (!B \land !A)$}
\DisplayProof
\]
Tingginya ialah $\pheight{\delta_4}+1=3$. Inferensi \Intro{\lif}
melepaskan semua kemunculan asumsi berbentuk $!A \land !B$ yang
berlabel~$x$, yakni kedua asumsi yang sebelumnya terbuka dalam
subbukti langsung. Karena keduanya merupakan satu-satunya asumsi
dari~$\delta_5$, bukti itu tidak memiliki asumsi terbuka, sehingga
merupakan !!a{proof} \emph{atas} !!{formula} akhirnya
$(!A \land !B) \lif (!B \land !A)$.
\end{ex}

\end{document}
